\documentclass[a4paper,12pt, reqno]{amsart}

\usepackage[svgnames]{xcolor}
%% Geometry adjustment 
%\usepackage[verbose]{geometry}
%\geometry{
%  paperheight=242mm,
%  paperwidth=165mm,
%  left=22.5mm,
%  right=22.5mm,
%  top=22.5mm,
%  bottom=22.5mm
%}
%\geometry{twoside=false,textwidth=380pt,textheight=560pt}
%\geometry{a5paper, hmargin=1cm, tmargin=1.3cm, bmargin=1.1cm}
\usepackage[centering, margin = 1in]{geometry}
%\setlength{\parskip}{2pt}

\usepackage[bbgreekl]{mathbbol}

\usepackage{amsmath, amssymb, amsthm, enumerate, url}
\usepackage{euscript}
\usepackage{amsfonts}
\usepackage{enumitem}
\usepackage{comment}
\usepackage{tikz-cd}
\usepackage{quiver}
\usepackage{mathrsfs}
\usepackage{mathtools}
\usepackage{graphicx}
\usepackage{thmtools}
\usepackage{braket}
\renewcommand{\Set}{\set}
\let\Set\relax

\usepackage{stmaryrd}
\SetSymbolFont{stmry}{bold}{U}{stmry}{m}{n}

%% bibilography management (old)
% \usepackage[style=alphabetic]{biblatex}
% \addbibresource{blibli.bib}

%% bibilography management (new)
% \usepackage[giveninits,style=alphabetic,useprefix=true,maxnames=9,maxalphanames=9,backend=biber,doi=false,isbn=false,useprefix=true]{biblatex}
% \addbibresource{filterization.bib}
% \usepackage[simplepostnote]{lumsdaine-biblatex-misc}

%% hyperref (old)
%\usepackage[colorlinks,citecolor=DarkRed,linkcolor=blue,urlcolor=blue]{hyperref}
%\usepackage{cleveref}

%% hyperref and friends (new)
\usepackage[svgnames]{xcolor}
\definecolor{darkgreen}{rgb}{0,0.45,0}
\definecolor{darkred}{rgb}{0.75,0,0}
\definecolor{darkblue}{rgb}{0,0,0.6}
\usepackage{hyperref}
\usepackage{cleveref}
\hypersetup{
	colorlinks = true,
  allcolors = NavyBlue,
}

\usepackage{soul}
\setul{3pt}{.4pt}
%\usepackage{pxfonts}
\usepackage{dsfont}

\usetikzlibrary{arrows.meta}
\newcommand{\ultrato}{\mathrel{\tikz[>={To}, line cap=round, anchor=base] \draw[-<] (0,0) -- (3.375mm,0);}}
\newcommand{\longultrato}{\mathrel{\tikz[>={To}, line cap=round, anchor=base] \draw[-<] (-13mm,0) -- (5mm,0);}}

%UA-command
\newcommand{\ult}{\mathrel{\tikz [line width=.11ex, double distance=.33ex]
    \draw[>-]
    (0.3,0) -- (0,0);}
}

%%%%%%%%%%%%%%%%%%%%%%%%%%%%%%%%%%%%%%%%%%%%%%%%%%%%
% Peter's ultra arrows

\usepackage{tikz}
\usepackage{tikz-cd}

%% Library imports
\usetikzlibrary{calc}
\usetikzlibrary{fit}
\usetikzlibrary{shapes}
\usetikzlibrary{arrows}
\usetikzlibrary{decorations.markings}
\usetikzlibrary{decorations.pathmorphing}
\usetikzlibrary{positioning}
\usetikzlibrary{intersections}

%%%% Set default styles for tikz-cd

\tikzset{
  commutative diagrams/arrow style=tikz,
%  commutative diagrams/diagrams={row sep=large},
}

%%%% Commands for using tikz-cd style in \tikzpicture 

\tikzset{cd-style/.style={commutative diagrams/every diagram}}
\tikzset{cd-arrow-style/.style={commutative diagrams/.cd, every arrow, every label}}
\newcommand{\arr}[1][]{\draw[cd-arrow-style,#1]}

% clone of how tikz-cd places arrow labels
% for use in tikzpicture environment as e.g. “\arr (A) to[label={f}{}] (B);” 
\makeatletter
\tikzset{
  label/.style n args={2}{
    edge node={node [
      execute at begin node=\iftikzcd@mathmode$\fi,%$
      execute at end node=\iftikzcd@mathmode$\fi,%$
      /tikz/commutative diagrams/.cd,every label,
      #2
      ] {#1}}
  }
}
\makeatother


%% Custom arrow tips & aliases

\tikzset{fibtip/.tip={Triangle[open,angle=60:4.5pt]}}
\tikzset{tfibtip/.tip={Bar[sep]Triangle[open,angle=45:4pt]}} % Based on http://tex.stackexchange.com/a/150213

\pgfarrowsdeclaredouble{fibbtip}{fibbtip}{fibtip}{.fibtip}

\tikzset{ulttip/.tip={Computer Modern Rightarrow[reversed]}}
% see PGF manual §16.5 for other arrow tip types 

\pgfarrowsdeclarealias{c}{c}{left hook}{right hook}
\pgfarrowsdeclarealias{c'}{c'}{right hook}{left hook}

%% Semantic arrow styles

% tip/tail styles: work in inline arrows, tikzpictures, and tikz-cd diagrams
\tikzset{fib/.code={\pgfsetarrowsend{fibtip}}}
\tikzset{fibb/.code={\pgfsetarrowsend{fibbtip}}}
\tikzset{tfib/.code={\pgfsetarrowsend{tfibtip}}}
\tikzset{inj/.code={\pgfsetarrowsstart{c}}}
\tikzset{inj'/.code={\pgfsetarrowsstart{c'}}}
\tikzset{cover/.style={->>}}
\tikzset{tcof/.style={tail}}
\tikzset{zigzag/.style={commutative diagrams/rightsquigarrow}}
\tikzset{ulttip/.code={\pgfsetarrowsend{ulttip}}}

% “label”-based styles: use in tikzpictures as “\arr (A) to[sub={\sigma}] (B)”
\tikzset{suplabel/.style={label={#1}{auto=left,pos=1}}}
\tikzset{sublabel/.style={label={#1}{auto=right,pos=1}}}
\tikzset{sub/.style={label={#1}{auto=right,pos=1}}}

% tikzcd-label-syntax-based styles: use in tikzcd diagrams
\tikzset{lw/.style={"lw"{#1}},lw/.default={very near end}}
\tikzset{lw'/.style={"lw"'{#1}},lw'/.default={very near end}}
\tikzset{weq/.style={"\sim"}}
\tikzset{weq'/.style={"\sim"'}}
\tikzset{ult/.style={ulttip,"#1"'{very near end}}}
\tikzset{ult'/.style={ulttip,"#1"{very near end}}}

%% Inline arrows

% TikZ inline arrows, adapted from http://tex.stackexchange.com/a/116324 and http://tex.stackexchange.com/a/188351

\newcommand{\generalto}[2]{ \mathrel{\mkern-1mu
  \tikz[baseline={([yshift=-0.58ex]a.south)}]{%
    \node[minimum width=1em,align=center,inner xsep=0.5ex,inner ysep=0.25ex] (a) {$\scriptstyle #2$};
    \draw[cd-arrow-style,#1] (a.south west) -- (a.south east);}
 \mkern-1mu}}

\newcommand{\generalfrom}[2]{ \mathrel{\mkern-1mu
  \tikz[baseline={([yshift=-0.58ex]a.south)}]{%
    \node[minimum width=1em,align=center,inner xsep=0.5ex,inner ysep=0.25ex] (a) {$\scriptstyle #2$};
    \draw[cd-arrow-style,#1] (a.south east) -- (a.south west);}
  \mkern-1mu}}

\renewcommand{\to}[1][]{ \generalto{->}{#1} }
\newcommand{\from}[1][]{ \generalfrom{->}{#1} }
\newcommand{\fibto}[1][]{ \generalto{fib}{#1} }
\newcommand{\fibbto}[1][]{ \generalto{fibb}{#1} }
\newcommand{\extto}[1][]{ \generalto{>->}{#1} }
\newcommand{\injto}[1][]{ \generalto{->,inj}{#1} }
\newcommand{\lwto}[1][]{ \generalto{->}{#1}_{\lw} }
\newcommand{\coverto}[1][]{ \generalto{->>}{#1\ \,} }
\newcommand{\vuto}[2][]{ \mathrel{{\generalto{ulttip}{#1}}\mkern2mu_{{#2}}\mkern-2mu}}
\newcommand{\vuleq}[2][]{ \mathrel{\preccurlyeq_{{#2}}^{{#1}}\mkern-2mu}}

%%%%%%%%%%%%%%%%%%%%%%%%%%%%%%%%%%%%%%%%%%%%%%%%%%%%%%%%%%%%%%%%%%%%%


% Errol: J'ai remplacé ça, normalement c'est mieux comme ça pour les références.
\setlist[enumerate]{label={(\roman*)}, ref={(\roman*)}} 
% \renewcommand{\labelitemi}{\boldmath$\cdot$}
% \renewcommand{\theenumi}{\roman{enumi}}

%\swapnumbers
\theoremstyle{plain}
\newtheorem{Thm}{Theorem}[section]
\newtheorem*{nThm}{Theorem}
\newtheorem{prop}[Thm]{Proposition}
\newtheorem{cor}[Thm]{Corollary}
\newtheorem{lem}[Thm]{Lemma}
\theoremstyle{definition}
\newtheorem{defn}[Thm]{Definition}
\newtheorem{para}[Thm]{}
\newtheorem{cons}[Thm]{Construction}
\newtheorem{nota}[Thm]{Notation}
\newtheorem*{nnota}{Notation}
\newtheorem{ex}[Thm]{Example}
\newtheorem{rmk}[Thm]{Remark}
\newtheorem*{nrmk}{Remark}
\newtheorem{claim}[Thm]{Claim}
\newtheorem*{nclaim}{Claim}
\newtheorem{question}[Thm]{Question}

\newcommand\pref[1]{\ref{#1}}

\DeclareMathOperator{\dom}{dom}
\DeclareMathOperator{\cod}{cod}
\DeclareMathOperator{\coeq}{coeq}
\DeclareMathOperator*{\Nat}{Nat}
\DeclareMathOperator*{\colim}{colim}
\DeclareMathOperator{\cl}{\mathsf{cl}}
\DeclareMathOperator{\subcl}{\mathsf{scl}}
 
% Arrows
\newcommand{\mono}{\hookrightarrow}
\newcommand{\epi}{\twoheadrightarrow}


\newcommand{\C}{\mathscr{C}}
\newcommand{\E}{\mathscr{E}}
\newcommand{\F}{\mathscr{F}}
\newcommand{\B}{\mathscr{B}}
\newcommand{\G}{\mathscr{G}}
\newcommand{\R}{\mathscr{R}}
\newcommand{\I}{\mathscr{I}}
\newcommand{\D}{\mathscr{D}}
\newcommand{\A}{\mathscr{A}}
\newcommand{\T}{\mathscr{T}}
\newcommand{\M}{\mathscr{M}}
\newcommand{\N}{\mathscr{N}}
\newcommand{\X}{\mathscr{X}}
\newcommand{\K}{\mathscr{K}}
\newcommand{\U}{\mathscr{U}}
\newcommand{\Z}{\mathscr{Z}}
\newcommand{\Y}{\mathscr{Y}}

\newcommand{\eC}{\EuScript{C}}
\newcommand{\eP}{\EuScript{P}}
\newcommand{\eO}{\EuScript{O}}
\newcommand{\eE}{\EuScript{E}}
\newcommand{\eF}{\EuScript{F}}
\newcommand{\eG}{\EuScript{G}}
\newcommand{\eB}{\EuScript{B}}
\newcommand{\eI}{\EuScript{I}}
\newcommand{\eD}{\EuScript{D}}
\newcommand{\eA}{\EuScript{A}}
\newcommand{\eT}{\EuScript{T}}
\newcommand{\eM}{\EuScript{M}}
\newcommand{\eZ}{\EuScript{Z}}
\newcommand{\eX}{\EuScript{X}}
\newcommand{\eY}{\EuScript{Y}}
\newcommand{\dsA}{\mathds{A}}
\newcommand{\dsD}{\mathds{D}}

\newcommand{\ptuf}{1}
\newcommand{\bb}{\bbbeta}
\newcommand{\En}{\mathrm{En}}
\newcommand{\Sh}{\mathsf{Sh}}
\newcommand{\sh}{\mathsf{sh}}
\newcommand{\vSh}{\mathfrak{sh}}
\newcommand{\Sub}{\mathsf{Sub}}
\newcommand{\Desc}{\mathsf{Desc}}
\newcommand{\Topos}{\mathsf{Topos}}
\newcommand{\Toposwep}{\mathsf{Topos}_{\mathrm{wep}}}
\newcommand{\CohTop}{\mathsf{CohTop}}
\newcommand{\Psh}{\mathsf{P}\Sh}
\newcommand{\vUlt}{\mathsf{vUlt}}
\newcommand{\vUltb}{\mathsf{vUlt}_{\mathrm{bounded}}}
\newcommand{\UMat}{\mathsf{UMat}}
\newcommand{\op}{\mathrm{op}}
\newcommand{\amp}{\mathsf{amp}}
\newcommand{\In}{\mathsf{in}}
\newcommand{\Cat}{\mathsf{Cat}}
\newcommand{\Pos}{\mathsf{Pos}}
\newcommand{\CAT}{\mathsf{CAT}}
\newcommand{\TopSp}{\mathsf{TopSp}}
\newcommand{\Ult}{\mathsf{Ult}}
\newcommand{\UltL}{\Ult^{\mathsf{L}}}
\newcommand{\pt}{\mathsf{pt}}
\newcommand{\vpt}{\mathfrak{pt}}
\newcommand{\id}{\mathsf{i}}
\newcommand{\diag}{\mathsf{d}}
\newcommand{\ev}{\mathsf{ev}}
\newcommand{\Set}{\mathsf{Set}}
\newcommand{\SET}{\mathsf{SET}}
\newcommand{\SetO}{\mathsf{Set}[\mathbb{O}]}
%\newcommand{\UF}{\EuScript{U}\EuScript{F}}
\newcommand{\UF}{\mathds{U}}
\newcommand{\Pretop}{\mathsf{Pretop}}
\newcommand{\Latt}{\mathsf{Latt}}
\newcommand{\Loc}{\mathsf{Loc}}
\newcommand{\Ind}{\mathsf{Ind}}
\newcommand{\FC}{\mathsf{FC}}
\newcommand{\Open}{\mathcal{O}}
\newcommand{\Ob}{\mathsf{Ob}}
\newcommand{\Hom}{\mathsf{Hom}}
\newcommand{\Mod}{\mathsf{Mod}}
\newcommand{\Alex}{\mathsf{Alex}}
\newcommand{\res}[1]{{\upharpoonright{#1}}}
\newcommand{\ie}{i.e., }
\newcommand{\eg}{e.g., }
\newcommand{\paronto}{\rightharpoonup\mathrel{\mspace{-15mu}}\rightharpoonup}
\newcommand{\li}{\mathsf{h}}
\newcommand{\lii}{\mathsf{k}}
\newcommand{\ty}{\mathsf{ty}}
\newcommand{\topo}{\mathsf{topo}}
%\usepackage{MnSymbol}
\newcommand{\vucat}{\mathsf{vucat}}
\newcommand{\vuCat}{\mathsf{vuCat}}
\newcommand{\vuCattaut}{\mathsf{vuCat_{taut}}}
\newcommand{\vuCAT}{\mathsf{vuCAT}}
\newcommand{\vuPos}{\mathsf{vuPos}}
\newcommand{\vuPOS}{\mathsf{vuPOS}}
\newcommand{\vuPostaut}{\mathsf{vuPos_{taut}}}
\newcommand{\vuPOStaut}{\mathsf{vuPOS_{taut}}}
\newcommand{\vupt}{\pt}

\newcommand{\sem}[1]{\llbracket{#1}\rrbracket}

\newcommand{\mrg}{\curlyvee}
\newcommand{\bigmrg}{\bigcurlyvee}

%\usepackage{bbm}
\newcommand{\Tref}{\rho_\mathrm{T}}
\newcommand{\Lref}{\rho_\mathrm{L}}
\newcommand{\Tinj}{\iota_\mathrm{T}}
\newcommand{\Linj}{\iota_\mathrm{L}}
\newcommand{\Et}{\mathrm{Et}}

\newcommand{\Ext}{\mathsf{Ext}}
\newcommand{\onespace}{\mathsf{1space}}
\newcommand{\Onespace}{\mathsf{1Space}}
\newcommand{\FF}{\mathds{F}}
\newcommand{\vf}{\mathsf{vf}}
\newcommand{\vfcat}{\mathsf{vfcat}}
\newcommand{\vfCat}{\mathsf{vfCat}}
\newcommand{\vfCattaut}{\mathsf{vfCat_{taut}}}
\newcommand{\vfCAT}{\mathsf{vfCAT}}
\newcommand{\vfPos}{\mathsf{vfPos}}
\newcommand{\vfPOS}{\mathsf{vfPOS}}
\newcommand{\vfPostaut}{\mathsf{vfPos_{taut}}}
\newcommand{\vfPOStaut}{\mathsf{vfPOS_{taut}}}
\newcommand{\vfpt}{\pt}

\newcommand{\Fam}{\mathsf{Fam}}
\newcommand{\Conn}{\mathsf{Conn}}
\newcommand{\Comp}{\mathsf{Comp}}
\DeclareFontFamily{U}{dmjhira}{}
\DeclareFontShape{U}{dmjhira}{m}{n}{ <-> dmjhira }{}
\DeclareRobustCommand{\yo}{\text{\usefont{U}{dmjhira}{m}{n}\symbol{"48}}}

\newcommand{\colorier}[1]{\color{Grey}{#1}\color{black}}


%\usepackage[inline]{showlabels} % comment to hide labels


%%% Comments 

\newcommand{\todo}[1]{{\color{darkred} todo: #1}}
\newcommand{\addref}{{\color{darkred} add reference }}
\newcommand{\citneed}{{\color{darkred} citation needed }}
\newcommand{\wip}{{\color{darkred} wip }}

\newcommand{\note}[1]{\textcolor{Blue}{#1}}
\newcommand{\UT}[1]{\textcolor{PaleVioletRed}{UT: #1}}
\newcommand{\gabrielnote}[1]{\textcolor{Teal}{GS: #1}}

\newcommand{\nocomments}{
    \renewcommand{\todo}[1]{}
    \renewcommand{\note}[1]{}
    \renewcommand{\gabrielnote}[1]{}
    \renewcommand{\errolnote}[1]{}
}

\newcommand{\emaillink}[1]{\href{mailto:#1}{\sf #1}}

\newcommand{\red}[1]{\textcolor{red}{#1}}
